% Part: normal-modal-logic
% Chapter: tableaux
% Section: introduction

\documentclass[../../../include/open-logic-section]{subfiles}

\begin{document}

\olfileid[pt]{nml}{tab}{int}

\olsection{Introdução}

Os !!^{tableau}s são certas árvores (ramificadas para baixo) de
!!{signed formula}s, O.s.\@, pares formados por um sinal de valor de
verdade ($\True$ ou $\False$) e !!a{sentence}
\[
\sFmla{\True}{!A} \text{ ou } \sFmla{\False}{!A}.
\]
!!^a{tableau} começa com certo número de \emph{hipóteses}. Cada
!!{signed formula} posterior é gerada pela aplicação de uma das regras
de inferência. Algumas regras de inferência acrescentam uma ou mais
!!{signed formula}s a uma extremidade da árvore; outras acrescentam duas
novas extremidades, produzindo dois ramos. As regras produzem !!{signed
formula}s nas quais a !!{formula} é menos complexa que a da !!{signed
formula} à qual a regra foi aplicada. Quando um ramo contém tanto
$\sFmla{\True}{!A}$ quanto $\sFmla{\False}{!A}$, dizemos que o ramo é
\emph{fechado}. Se todos os ramos de !!a{tableau} são fechados, o
!!{tableau} inteiro é fechado. Um !!{tableau} fechado constitui
!!a{derivation} que mostra que o conjunto das !!{signed formula}s usadas
para iniciar o !!{tableau} é insatisfatível. Isso pode ser usado para
definir uma relação $\Proves$: $\Gamma \Proves !A$ sse existe algum
conjunto finito~$\Gamma_0 = \{!B_1, \dots, !B_n\} \subseteq \Gamma$ tal
que há um !!{tableau} fechado para as hipóteses
\[
\{\sFmla{\False}{!A}, \sFmla{\True}{!B_1}, \dots, \sFmla{\True}{!B_n}\}.
\]

Para as lógicas modais, precisamos estender a noção de !!{signed
formula} e acrescentar regras que tratem de~\iftag{prvBox}{$\Box$
\iftag{prvDiamond}{ e $\Diamond$}}{$\Diamond$}. Além de um sinal
($\True$ ou $\False$), as !!{formula}s em !!{tableau}s modais também
têm \emph{prefixos}~$\sigma$. Os prefixos são sequências não vazias de
inteiros positivos, O.s.\@, $\sigma \in (\PosInt)^* \setminus
\{\emptyseq\}$. Escrevemos tais prefixos sem os delimitadores
$\tuple{\ }$ e separamos os !!{element}s individuais por~$.$, em vez
de~$,$. Se $\sigma$ é um prefixo, então $\sigma.n$ é $\sigma
\concat \tuple{n}$; p.ex., se $\sigma = 1.2.1$, então $\sigma.3$
é $1.2.1.3$. Assim, por exemplo,
\[
\sFmla{\True}{\Box !A \lif !A}[1.2]
\]
é uma \emph{!!{signed formula} prefixada} (ou, abreviadamente, apenas uma
\emph{!!{formula} prefixada}).

Intuitivamente, o prefixo nomeia um mundo de um modelo que poderia
satisfazer as !!{formula}s de um ramo de !!a{tableau}; se $\sigma$
nomeia algum mundo, então $\sigma.n$ nomeia um mundo acessível a partir
do mundo nomeado por~$\sigma$.

\end{document}
